\documentclass[12pt, openright, a4paper, twoside, english, brazil]{article}
\usepackage[top=3cm, bottom=2cm, left=3cm, right=2cm]{geometry}

\usepackage{mathtools}
\usepackage{amsfonts}
\usepackage{amssymb}
\usepackage{amsmath}
\usepackage{wasysym}
\usepackage[utf8]{inputenc}
\usepackage[portuguese]{babel}
\usepackage[T1]{fontenc}

\usepackage{parskip}
\usepackage{xcolor}
\usepackage{cancel}
\usepackage{enumerate}

\usepackage{libertine}
\usepackage{newtxmath}

\newenvironment{qitem}{
  \color{black}
\item}{\vspace{10pt}}

\newenvironment{qmitem}{
  \color{black}
\item
\begin{enumerate}[\bfseries a)]}{
\end{enumerate}\vspace{10pt}}

% -----------------------------------------------------------------------------

\title{Soluções - FME, Volume 06}
\author{Capítulo I - Exercícios}
\date{}

\begin{document}

\maketitle

\begin{enumerate}[\bfseries 1.]
  \begin{qmitem} % QUESTION 1
  \item $(6+7i)(1+i)\color{red}=(6,7)(1,1) = (6-7, 6+7) = -1+13i$
  \item $(5+4i)(1-i)+(2+i)i\color{red}=(5,4)(1,-1) + (2,1)(0,1) =
    (5+4+0-1, -5+4+2-0) = 8+i$
  \item $(1+2i)^2-(3+4i)\color{red}=(1,2)^2 - (3,4) = (1-2^2, 4) - (3,4) = -6$
  \end{qmitem}

  \begin{qmitem} % QUESTION 2
  \item$(3+2i) + (2-5i)\color{red} = 5-3i$
  \item$(5-2i) - (2+8i)\color{red} = 3-10i$
  \item$(1+i) + (1-i) - 2i\color{red} = 2-2i$
  \item$(6+7i) - (4+2i) + (1-10i)\color{red} = 3-5i$
  \end{qmitem}

  \begin{qmitem} % QUESTION 3
  \item $(2-3i)(1+5i) =\color{red} (2+15, 10-3) = 17+7i$
  \item $(1+2i)(2+i) =\color{red} (2-2, 1+4) = 5i$
  \item $(4-3i)(5-i)(1+i) =\color{red} (20-3, -4-15)(1,1) = (17+19,
    17-19) = 36-2i$
  \item $(7+2i)(7-2i) =\color{red} (49+4, -14+14) = 53$
  \end{qmitem}

  \begin{qitem} % QUESTION 4
    $u=4+3i, v=5-2i, uv=?$

    $\color{red} uv=(4,3)(5,-2) = (20+6, -8+15) = 26+7i$
  \end{qitem}

  \begin{qmitem} % QUESTION 5
  \item $(3+2i)^2 \color{red}=(3,2)^2 = (3^2-2^2, 2 \cdot 6) = 5+12i$
  \item $(5-i)^2 \color{red}=(5,-1)^2 = (5^2-(-1)^2, -10) = 24-10i$
  \item $(1+i)^3 \color{red}=(1,1)^3 = (1^2-1^2, 2)(1,1) = (0-2, 0+2) = -2+2i$
  \end{qmitem}

  \begin{qitem} % QUESTION 6
    $f(z) = z^2 - z + 1, f(1-i) = ?$

    $\color{red}f(1-i)= (1,-1)^2 - (1,-1) + (1,0) = (0,-2) + (0,1) = -i$
  \end{qitem}

  \begin{qitem} % QUESTION 7
    $f(z) = z^4 + iz^3 - (1+2i)z^2 + 3z + 1 + 3i, f(1+i) = ?$

    $\color{red}f(1+i)=(1,1)^2 = (0,2)$

    $\color{red}= (0,2)^2 + (0,1)(0,2)(1,1) - (1,2)(0,2) + (3,3) + (1,3)$

    $\color{red}= (-4,0) + (-2,0)(1,1) - (-4,2) + (4,6)$

    $\color{red}= 2[(-2,0) + (-1,-1) - (-2,1) + (2,3)]$

    $\color{red}= 2(1,1) = 2(1+i)$
  \end{qitem}

  \begin{qitem} % QUESTION 8
    Prove que $(1+i)^2 = 2i$ e coloque na forma algébrica $z =
    \frac{(1+i)^{80} - (1+i)^{82}}{i^{96}}$

    \begin{itemize}
      \item[1.] $\color{red} (1,1)^2 = (1-1, 1+1) = 2i$
      \item[2.] $\displaystyle \color{red} z = \frac{(2i)^{40} -
        (2i)^{41}}{(i^2)^{48}} =
        \frac{2^{40} - 2^{41} \cdot i}{1} = 2^{40} - 2^{41}i$
    \end{itemize}
  \end{qitem}

  \begin{qmitem} % QUESTION 9
  \item $\color{red} i^{76} = (i^2)^{38} = 1$
  \item $\color{red} i^{110} = (i^2)^{55} = -1$
  \item $\color{red} i^{97} = (i^2)^{48} \cdot i = i$
  \item $\color{red} i^{503} = (i^2)^{251} \cdot i = -i$
  \end{qmitem}

  \begin{qitem} % QUESTION 10
    Prove que $(1-i)^2 = -2i$ e calcule $(1-i)^{96} + (1-i)^{97}$

    $\color{red} (1,-1)^2 = (1-(-1)^2, -1-1) = -2i$

    $\color{red} (1,-1)^{96} + (1,-1)^{97} = (-2i)^{48} +
    (-2i)^{48}(1,-1) = 2^{48} + 2^{48} - 2^{48}i = 2^{49} - 2^{48}i$
  \end{qitem}

  \begin{qitem} % QUESTION 11
    Se $i^2 = -1$, calcule $(1+i)^{12} - (1-i)^{12}$

    $\color{red}= (2i)^6 - (-2i)^6 = -2^6 + 2^6 = 0$
  \end{qitem}

  \begin{qitem} % QUESTION 12
    Simplifique $\displaystyle
    \frac{(2+i)^{101}(2-i)^{50}}{(-2-i)^{100}(i-2)^{49}}$

    $=\color{red}\displaystyle
    \frac{(2,1)^{101}(2,-1)^{50}}{(-(2,1))^{100}(-(2,-1))^{49}} =
    \frac{(2,1)^{101}(2,-1)^{50}}{(2,1)^{100}(-1)(2,-1)^{49}} =
    -(2,1)(2,-1) = -(4+1, -2+2) = -5$
  \end{qitem}

  \begin{qitem} % QUESTION 13
    A igualdade $(1+i)^n = (1-i)^n$ é válida para os números
    naturais divisíveis por qual natural?

    \color{red}
    \begin{itemize}
      \item Para $n = 2k, k \in \mathbb{N} \Rightarrow ((1+i)^2)^k =
        ((1-i)^2)^k \Rightarrow (2i)^k = (-2i)^k \Rightarrow 2^k i^k =
        (-1)^k 2^k i^k$

        $\Rightarrow (-1)^k = 1 \Rightarrow k = 2k'$

        $\Rightarrow n = 4k'', k'' \in \mathbb{N}$

      \item Para $n = 2k+1, k \in \mathbb{N} \Rightarrow ((1+i)^2)^k(1+i) =
        ((1-i)^2)^k(1-i) \Rightarrow (2i)^k(1+i) = (-2i)^k(1-i)$

        $\displaystyle \Rightarrow \frac{(1+i)(1+i)}{(1-i)(1+i)} = (-1)^k
        \Rightarrow (-1)^k = i \implies \bot$
    \end{itemize}

    $\therefore n\text{ deve ser divisível por 4.}$
  \end{qitem}

  \begin{qmitem} % QUESTION 14
  \item $\color{red}2+3yi = x+9i \quad x=2, 3y=9 \Rightarrow
    \begin{cases} x=2 \\ y=3
    \end{cases}$

  \item $\color{red}(x+yi)(3+4i) = 7+26i \quad (3x-4y, 4x+3y) =
    (7,26) \Rightarrow
    \begin{cases} 3x-4y=7 \\ 4x+3y=26
    \end{cases}$

    $\color{red} \cancel{3x} - 4y - \frac{3}{\cancel{4}}\cancel{4x} -
    \frac{3}{4}
    \cdot 3y = 7 - \frac{3}{4} \cdot 26 \Rightarrow -4y - \frac{9}{4}y =
    7 - \frac{3 \cdot 26}{4} \Rightarrow -y\left(\frac{4^2+3^2}{4}\right)
    = \frac{28-78}{4}$

    $\color{red} y = \frac{50}{25} = 2$

    $\color{red}\Rightarrow x = \frac{7 + 4 \cdot 2}{3} = \frac{15}{3} =
    5 \Rightarrow
    \begin{cases} x=5 \\ y=2
    \end{cases}$

  \item $\color{red} (x+yi)^2 = 4i \quad (x^2-y^2, 2xy) = (0,4) \Rightarrow
    \begin{cases} x^2=y^2 \Rightarrow x=y \\ xy=2
    \end{cases} \Rightarrow x^2=2 \Rightarrow
    \begin{cases} x=y=\sqrt{2} \\ \text{ou} \\ x=y=-\sqrt{2}
    \end{cases}$
  \end{qmitem}

  \begin{qmitem} % QUESTION 15
  \item $\color{red} 3+5ix = y-15i \Rightarrow
    \begin{cases} y=3 \\ 5x=-15
    \end{cases} \Rightarrow
    \begin{cases} x=-3 \\ y=3
    \end{cases}$

  \item $\color{red} (x+yi)(2+3i) = 1+8i \quad (2x-3y, 3x+2y) = (1,8)
    \Rightarrow
    \begin{cases} 2x-3y=1 \\ 3x+2y=8
    \end{cases}$

    $\color{red}\displaystyle \cancel{2x} - 3y -
    \frac{2}{\cancel{3}}\cancel{3x} -
    \frac{2}{3} \cdot 2y = 1 - \frac{2}{3} \cdot 8 \Rightarrow
    y\left(\frac{-9-4}{3}\right) = \frac{3-16}{3} \Rightarrow y =
    \frac{-13}{-13} = 1$

    $\color{red} \Rightarrow x = \frac{1+3}{2} \Rightarrow
    \begin{cases} x=2 \\ y=1
    \end{cases}$

  \item $\color{red} (3+yi) + (x-2i) = 7-5i \Rightarrow
    \begin{cases} 3+x=7 \\ y-2=-5
    \end{cases} \Rightarrow
    \begin{cases} x=4 \\ y=-3
    \end{cases}$

  \item $\color{red} (x+yi)^2 = 2i \Rightarrow (x^2-y^2, 2xy) = (0,2)
    \Rightarrow
    \begin{cases} x^2=y^2 \\ xy=1
    \end{cases} \Rightarrow
    \begin{cases} x=1 \text{ e } y=1 \\ \text{ou} \\ x=-1 \text{ e } y=-1
    \end{cases}$

  \item $\color{red} (2-x+3y) + 2yi = 0 \Rightarrow
    \begin{cases} 2-x+3y=0 \\ 2y=0
    \end{cases} \Rightarrow
    \begin{cases} x=2 \\ y=0
    \end{cases}$

  \item $\color{red} (3-i)(x+yi) = 20 \quad (3x+y, 3y-x) = (20,0) \Rightarrow
    \begin{cases} 3x+y=20 \\ 3y=x
    \end{cases}$

    $\color{red}\displaystyle\Rightarrow 3x + \frac{x}{3} = 20
    \Rightarrow 9x+x = 3
    \cdot 20 \Rightarrow x=6 \Rightarrow
    \begin{cases} x=6 \\ y=2
    \end{cases}$
  \end{qmitem}

  \begin{qitem} % QUESTION 16
    $
    \begin{cases} x+yi=i \\ xi+y=2i-1
    \end{cases}$

    $\color{red}x=i(1-y) \Rightarrow (1-y)i^2+y=2i-1 \Rightarrow
    y-1+y=2i-1 \Rightarrow y=i$

    $\color{red}\Rightarrow x+i^2=i \Rightarrow x=1+i \Rightarrow
    \begin{cases} x=1+i \\ y=i
    \end{cases}$
  \end{qitem}

  \begin{qitem} % QUESTION 17
    $(a,b)(c,d) = (ac-bd, ad+bc) = (x,0), x \in \mathbb{R}$

    $\color{red}\displaystyle\Rightarrow ad+bc=0 \Rightarrow ad=-bc \Rightarrow
    \frac{a}{c} = -\frac{b}{d}$
  \end{qitem}

  \begin{qitem} % QUESTION 18
    $(a,b)^2 \Rightarrow (a^2-b^2, 2ab)^2 = ((a^2-b^2)^2 - 4a^2b^2,
    4ab(a^2-b^2))$

    \color{red}
    $\Rightarrow 1. \ 4ab(a^2-b^2) = 0 \Rightarrow a = \pm b \text{ (}a
    \neq 0 \wedge b \neq 0\text{) ou } a=0 \text{ ou } b=0$

    $\Rightarrow 2. \ (a^2-b^2)^2 - 4a^2b^2 < 0 \Rightarrow
    (a^2-b^2)^2 < 4a^2b^2$
    \begin{itemize}
      \item[$\bullet$] $a=0, b \neq 0 \Rightarrow (-b^2)^2 < 0 \Rightarrow
        b^4 < 0 \Rightarrow \bot$
      \item[$\bullet$] $a \neq 0, b=0 \Rightarrow a^4 < 0 \Rightarrow \bot$
      \item[$\bullet$] $a = \pm b, a \neq 0, b \neq 0 \Rightarrow ((\pm
        b)^2 - b^2)^2 < 4(\pm b)^2 b^2 \Rightarrow 0 < 4b^4, \forall b \neq 0$
    \end{itemize}

    $\therefore
    \begin{cases} a = \pm b \\ ab \neq 0
    \end{cases}$
  \end{qitem}

  \begin{qmitem} % QUESTION 19
  \item $\displaystyle \frac{2}{i}\color{red}\cdot\frac{-i}{-i}=-2i$

  \item $\displaystyle
    \frac{3}{2+i}\color{red}\cdot\frac{2-i}{2-i}=\frac65 -\frac35i$

  \item $\displaystyle
    \frac{1+2i}{3-i}\color{red}\cdot\frac{3+i}{3+i}=\frac1{10} +\frac7{10}i$

  \item $\displaystyle
    \frac{i^9}{4-3i}\color{red}\cdot\frac{4+3i}{4+3i}=\frac{i^8\cdot
    i(4+3i)}{25}=\frac{-3}{25}+\frac{4}{25}i$
  \end{qmitem}

  \begin{qmitem} % QUESTION 20
  \item $\displaystyle \frac{1}{i}\color{red}\cdot\frac{i}{i}=-i$

  \item $\displaystyle
    \frac{1}{1+i}\color{red}\cdot\frac{1-i}{1-i}=\frac12-\frac12i$

  \item $\displaystyle
    \frac{3+4i}{2-i}\color{red}\cdot\frac{2+i}{2+i}=\frac23+\frac{11}5i$

  \item $\displaystyle
    \frac{1+i}{1-i}\color{red}\cdot\frac{1+i}{1+i}=\frac{2i}2=i$

  \item $\displaystyle
    \frac{i^{11}+2i^{13}}{i^{18}-i^{37}}\color{red}=\frac{i^{11}(1+2i^2)}{i^{18}(1-i^19)}=\frac{-1}{i^4i^2i(1-i^{16}i^2i)}=\frac{1}{i(1+i)}=\frac{1}{-1+i}\cdot\frac{-1+i}{-1+i}=-\frac12-\frac12i$

  \item $\displaystyle
    \frac{1-3i}{3-i}\color{red}\cdot\frac{3+i}{3+i}=\frac35-\frac45i$

  \item $\displaystyle
    \frac{i^3-i^2+i^{17}-i^{35}}{i^{16}-i^{13}+i^{30}}\color{red}=\frac{i^2(i-1+i^{12}i^3-i^{32}i)}{i^{13}(i^3-1+i^{17})}=\frac{i-1-i-i}{i^3(-i-1+i^{16}i)}=\frac{-(1+i)}{-i(-1)}\cdot\frac{-i}{-i}=-1+i$

  \item $\displaystyle
    \frac{1+i}{(1-i)^2}\color{red}=\frac{1+i}{-2i}\cdot\frac{2i}{2i}=\frac{-2+2i}{4}=-\frac12+\frac12i$
  \end{qmitem}

  \begin{qitem} % QUESTION 21
    $\displaystyle
    \frac{1+3i}{2-i}\color{red}\cdot\frac{2+i}{2+i}=-\frac15-\frac75i$
  \end{qitem}

  \begin{qitem} % QUESTION 22
    $\displaystyle
    \frac{1+i}{i}\color{red}\cdot\frac{-i}{-i}=1-i\implies 1+i$
  \end{qitem}

  \begin{qitem} % QUESTION 23
    $\displaystyle
    z=\frac{1-i}{1+i}\color{red}\cdot\frac{-2i}{2}=-i\implies
    z^{-1}=\frac{1}{-i}\cdot\frac{i}{i}=i\implies \overline{z^{-1}}=-i$
  \end{qitem}

  \begin{qitem} % QUESTION 24
    $u^2-v^2=6,\bar u+\bar v=1-i,u-v=?$

    \color{red}
    $\displaystyle (u-v)(u+v)=6\implies u-v=\frac{6}{u+v}$

    $\bar u+\bar v=1-i\implies \overline{\bar u+\bar
    v}=\overline{1-i}\implies u+v=1+i$

    $\displaystyle \therefore
    u-v=\frac{6}{u+v}=\frac{6}{1+i}\cdot\frac{1-i}{1-i}=3-3i$
  \end{qitem}

  \begin{qitem} % QUESTION 25
    \text{Provavelmente tem uma informação errada na questão} $u=1$
  \end{qitem}

  \begin{qitem} % QUESTION 26
    $\displaystyle\left(\frac{1+i}{1-i}\right)^n, n\in\mathbb{Z}$

    \color{red}

    $\displaystyle
    \left(\frac{1+i}{1-i}\cdot\frac{1+i}{1+i}\right)^n=\left(\frac{2i}{2}\right)^n=i^n\therefore
    1,i,-1,-i$
  \end{qitem}

  \begin{qitem} % QUESTION 27
    $\displaystyle \frac{a+bi}{c+di}, c+di\ne 0$

    \color{red}
    $\displaystyle
    \frac{a+bi}{c+di}\cdot\frac{c-di}{c-di}=\frac{ac-bd+i(ad+bc)}{c^2+d^2}$

    $\implies  ax-bd=0$

    $\implies  ad+bc=0$
  \end{qitem}

\end{enumerate}

\end{document}
